\section{Formal Mathematical Proofs}

\begin{theorem}[Selective Noise Shredding Ratio Theorem]
Let $\theta_{\text{noise}}$ satisfy $|\theta_{\text{noise}}| \le 1.0$ and $\theta_{\text{macro}}$ satisfy $|\theta_{\text{macro}}| > 1.0$. Under Lucas-Pell Hybrid Weight Decay, the decay rate ratio of noise to macro-structure parameters is exactly $\Phi \cdot \delta_{\text{Pell}} = \frac{1+\sqrt{5}}{2}(1+\sqrt{2}) \approx 3.90623$, shredding noise nearly $3.9\times$ faster.
\end{theorem}

\begin{proof}
Comparing the decay modifiers:
\begin{equation}
    \frac{\eta(\theta_{\text{noise}})}{\eta(\theta_{\text{macro}})} = \frac{\Phi}{\delta_{\text{Pell}}^{-1}} = \Phi \cdot \delta_{\text{Pell}}
\end{equation}
Substituting numerical values $\Phi \approx 1.61803399$ and $\delta_{\text{Pell}} \approx 2.41421356$:
\begin{equation}
    \Phi \cdot \delta_{\text{Pell}} = 1.61803399 \times 2.41421356 \approx 3.9062326
\end{equation}
Thus, chaotic background noise undergoes exponential decay at $3.906\times$ the rate of dominant structural weights, establishing selective noise shredding.
\end{proof}

\begin{theorem}[Macro-Structure Bounded Conservation Theorem]
For dominant weights $|\theta| > 1.0$, the parameter retention factor per step is $1 - \lambda_{\text{base}} \delta_{\text{Pell}}^{-1} \approx 1 - 0.414 \lambda_{\text{base}}$, preserving structural capacity compared to standard L2 decay ($1 - \lambda_{\text{base}}$).
\end{theorem}

\begin{proof}
Under standard L2 decay, the retention factor is $1 - \lambda_{\text{base}}$. Under Lucas-Pell decay for $|\theta| > 1.0$, the retention factor is $1 - \frac{\lambda_{\text{base}}}{1+\sqrt{2}} \approx 1 - 0.4142 \lambda_{\text{base}}$.
Since $0.4142 < 1.0$, $1 - 0.4142 \lambda_{\text{base}} > 1 - \lambda_{\text{base}}$, proving superior macro-structural parameter retention.
\end{proof}